\documentclass[11pt]{article}

\usepackage[margin=1in]{geometry}
\usepackage[T1]{fontenc}
\usepackage[utf8]{inputenc}
\usepackage{lmodern}
\usepackage{amsmath,amssymb,amsthm}
\usepackage{array}
\usepackage{booktabs}
\usepackage{graphicx}
\usepackage{subcaption}
\usepackage[sort,comma]{natbib}
\usepackage{hyperref}
\usepackage{xcolor}
\usepackage{enumitem}

\newtheorem{lemma}{Lemma}

\hypersetup{
  colorlinks=true,
  linkcolor=blue,
  citecolor=blue,
  urlcolor=blue
}

\graphicspath{{../results/benchmark_revision/figures/}}

\title{A General Model of Network Growth: Preferential Attachment with Capacity, Cost, and Connectivity}
\author{Nazanin Tajik \and Somwrita Sarkar \and Alireza Ermagun \and David Levinson}
\date{\today}

\begin{document}

\maketitle

\begin{abstract}
We propose a unifying, access-based model of network growth in which the probability that a new link attaches to an existing node is proportional to the marginal change in access, net of cost. Access is an auditable and measurable quantity in transport, communication, and social systems; costs capture distance, delay, or monetary expense; and capacity captures the fact that attractive nodes are not infinitely expandable. The resulting attachment kernel is tunable and nests Barab\'asi--Albert preferential attachment, spatial attachment, and capacity-limited growth as special cases. The contribution is not the invention of each ingredient in isolation, but their synthesis in one dyadic model and the systematic mapping of the regimes they generate. In a revision benchmark with deeper replications up to $N=1000$, finite capacity truncates hub growth: at $N=1000$, imposing $K=16$ reduces the mean maximum degree from $82.31$ in the Barab\'asi--Albert benchmark to $16.00$ in the capacity-constrained cases. Increasing distance deterrence from $\phi=0$ to $\phi=2$ reduces mean edge length from $0.526$ to $0.157$ and raises clustering from $0.007$ to $0.072$. Increasing multi-link entry from $\kappa=1$ to $\kappa=4$ raises the cyclomatic number from $6$ to $2991$ and reduces average shortest-path length from $8.19$ to $3.58$. A heterogeneous-capacity extension shows that dispersion in $K_i$ partially restores hub inequality: under lognormal capacities with mean near $16$, the mean maximum degree rises to $48.69$ and the degree Gini to $0.382$. Tail fits indicate that the capacity-constrained cases favor exponential tails, while the BA benchmark only clearly favors a power-law tail at the larger executed size. Within the explored benchmark envelope, no runs terminated early and all benchmark networks remained connected. The model therefore identifies when hub-dominated, localized, saturated, and meshed network forms emerge, and it generates falsifiable hypotheses for future empirical work.
\end{abstract}

\section{Introduction}

Preferential attachment remains one of the canonical explanations for the emergence of heavy-tailed degree distributions in growing networks \citep{barabasi1999emergence}. Yet many real networks, especially physical and social systems, do not evolve in purely topological space. Roads, rail systems, airline networks, pipelines, and social relations are embedded in geography, face finite capacity, and often allow entering nodes to form more than one link. In such systems, the attractiveness of a target depends not only on its current connectivity, but also on residual capacity and the dyadic cost of reaching it.

We posit that links form to increase \emph{marginal access}---the change in the ability of an origin to reach valued destinations---net of cost. This ``access first'' interpretation is empirically appealing because access can be measured, audited, and linked to infrastructure performance \citep{levinson2020towards,wu2020unifying,wang2022time}. It also provides a transparent micro-foundation for why attachment should depend jointly on connectivity, residual capacity, and distance or generalized cost.

The need for such a generalization is sharpened by the literature questioning universal scale-freeness. Critics have argued that many observed networks are not truly scale-free, and that power laws are rarer than often claimed \citep{holme2019rare,broido2019scale,clauset2009power,willinger2009mathematics}. More generally, the same empirical regularities that make preferential attachment attractive also motivate asking when heavy tails should break down. Capacity and distance are natural candidates for that breakdown.

\paragraph{Contributions.}
\begin{enumerate}[leftmargin=1.5em]
    \item We formalize a marginal-access growth rule and derive a compact attachment kernel that embeds degree effects, capacity slack, and spatial friction.
    \item We show explicitly how the model nests Barab\'asi--Albert preferential attachment as a limiting case and provide a dictionary linking parameter settings to well-known model families.
    \item We use simulation to map the comparative statics, interactions, and regime structure of the model, rather than merely illustrating that it can draw alternative network pictures.
    \item We generate falsifiable hypotheses for future empirical work on transport and infrastructure network growth, while deliberately keeping this paper as a theory-plus-simulation contribution.
\end{enumerate}

The paper is aimed at \emph{Network Science}. Its purpose is to show what the generalized model teaches that the equation alone does not: the shape of the outcome space, the conditions under which BA-like intuition survives, and the ways capacity, distance, and multi-link entry jointly bend growth toward localized, saturated, or meshed regimes.

\section{Related Work}

Preferential attachment explains broad classes of heavy-tailed networks \citep{barabasi1999emergence}, but spatial and geometry-aware variants show that distance deterrence affects clustering, path lengths, and tail exponents \citep{barthelemy2011spatial}. Capacity- or age-limited attachment can generate exponential or stretched-exponential cutoffs \citep{krapivsky2001organization,bianconi2001competition}. Fitness and hidden-metric models replace popularity with latent heterogeneity, while work on ``scale-free is rare'' emphasizes the need for more careful tail diagnostics and more modest claims of universality \citep{clauset2009power,holme2019rare,broido2019scale,voitalov2019scale}.

Our contribution is not to claim novelty for each ingredient in isolation. Spatial preferential attachment already exists; saturation and cutoff variants already exist. What is new here is the attempt to combine preferential attachment, capacity, dyadic cost, and multi-link entry in one access-based formulation, and then to map the resulting behavioral domains systematically. In that sense, the goal is synthesis plus regime identification rather than a standalone new kernel with no organizing interpretation.

The scale-free critique literature is especially important for framing the paper. A large body of work has shown that many empirical networks deviate from idealized power-law claims and that inference based on straight lines in log-log plots is inadequate \citep{newman2003structure,dorogovtsev2000generic,li2005towards,clauset2009power,golosovsky2017power,broido2019scale,smith2021scarcity}. That literature suggests that a useful growth model should explain not only why heavy tails can arise, but also why they may be truncated, localized, or absent.

\section{Model}
\label{sec:model}

\subsection{Access-based link formation}

Access is the ability to reach valued destinations \citep{levinson2020towards}. Let $G_t=(V_t,E_t)$ be the network at step $t$. For any origin $o$, define generalized access as
\begin{equation}
A_o(G_t)=\sum_{d\in V_t} g(O_d)\,f(C_{od}(G_t)),
\label{eq:access}
\end{equation}
where $O_d$ is destination mass, $g(\cdot)$ transforms opportunity, $C_{od}(G_t)$ is generalized cost over the current network, and $f(\cdot)$ is a decreasing impedance function.

When a new node considers attaching to an existing target $i$, the central idea is that attachment is proportional to the increase in access net of cost:
\begin{equation}
\Pi_{o\to i} \propto \Delta A_{o\to i},
\qquad
\Delta A_{o\to i}=A_o(G_t\cup (o,i)) - A_o(G_t).
\label{eq:deltaA}
\end{equation}
This statement is deliberately general. The compact kernel used in the simulations below is a reduced-form approximation to this access logic that makes the mechanism estimable and computationally tractable.

To connect Equation~\ref{eq:deltaA} to an estimable attachment rule, we assume that the ranking of candidate targets can be approximated by a separable local expansion in three monotone components: accumulated reach, residual slack, and dyadic cost. Let $s_i=\max(K_i-k_i,0)$ denote residual capacity. We then write
\begin{equation}
\Delta A_{n\to i} \approx \xi\,(k_i+\varepsilon)^{\alpha}\,s_i^{\beta}\,h(c_{ni}),
\label{eq:deltaA-approx}
\end{equation}
where $\xi$ is a proportionality constant and $h(\cdot)$ is decreasing in cost. Choosing the parsimonious power form $h(c_{ni})=(c_{ni}+\varepsilon)^{-\phi}$ yields the kernel used below. The point is not that Equation~\ref{eq:kernel} is a fully solved closed form for accessibility, but that it preserves the signs and elasticities implied by the marginal-access logic while remaining tractable for simulation and later estimation.

\subsection{Generalized attachment kernel}

At each time step a new node $n$ arrives, is placed in the unit square, and attempts to form $\kappa$ links to feasible existing nodes. Let $k_i$ denote the current degree of node $i$, let $K_i$ denote its capacity, and let $c_{ni}$ denote the dyadic cost between the arriving node $n$ and target $i$. We define the default attachment weight as
\begin{equation}
w_{ni} = (k_i+\varepsilon)^{\alpha}\,\max(K_i-k_i,0)^{\beta}\,(c_{ni}+\varepsilon)^{-\phi},
\label{eq:kernel}
\end{equation}
where $\alpha \ge 0$ captures economies of scale or preferential attachment, $\beta \ge 0$ captures saturation pressure, $\phi \ge 0$ captures distance deterrence, and $\varepsilon > 0$ is a small numerical constant.

The choice probability is then
\begin{equation}
\Pi_{ni}=\frac{w_{ni}}{\sum_{j\in\mathcal{F}_n}w_{nj}},
\qquad
\mathcal{F}_n=\{j:K_j>k_j\}.
\label{eq:probability}
\end{equation}
Neighbors are chosen sequentially without replacement within each arrival step, and degrees are updated after every realized edge. If fewer than $\kappa$ feasible targets remain, the node attaches to all feasible targets; if none remain, growth stops.

This kernel is consistent with the access interpretation. High-degree nodes proxy accumulated access advantages, residual capacity captures diminishing marginal value as a node fills, and dyadic cost discounts the value of distant or expensive links.

\subsection{PA as an access limit}

\begin{lemma}[Preferential attachment as a special case of access]
Let accessibility be given by Equation~\ref{eq:access}. Assume (i) identical destination mass, $g(O_d)\equiv 1$; (ii) a step impedance on topological distance, $f(C_{od})=\mathbf{1}\{\mathrm{dist}_G(o,d)\le 1\}$; (iii) identical generalized link costs or costs ignored at link choice; and (iv) a sparse, locally tree-like network so neighborhood overlap is negligible at the time of attachment. If a new node chooses target $i$ with probability proportional to the marginal access gain $\Delta A_{o\to i}$, then
\[
\Pi_{o\to i}\propto k_i,
\]
so the rule reduces to Barab\'asi--Albert preferential attachment.
\end{lemma}

\begin{proof}[Sketch]
Under assumptions (i) and (ii), access counts immediate neighbors. Attaching to target $i$ exposes the new node to $i$ and approximately to the $k_i$ neighbors of $i$, with overlap negligible by assumption (iv). Hence $\Delta A_{o\to i}\approx 1+k_i$, so after dropping the constant term one obtains $\Delta A_{o\to i}\propto k_i$.
\end{proof}

\subsection{Model dictionary}

Table~\ref{tab:dictionary} summarizes how prominent model families arise as limiting cases of Equation~\ref{eq:kernel}.

\begin{table}[htbp]
\centering
\caption{Model dictionary: limiting cases of the generalized attachment kernel.}
\label{tab:dictionary}
\begin{tabular}{@{}>{\raggedright\arraybackslash}p{3.0cm}>{\raggedright\arraybackslash}p{4.0cm}>{\raggedright\arraybackslash}p{6.0cm}@{}}
\toprule
Family & Parameters & Interpretation \\
\midrule
Barab\'asi--Albert (PA) & $\alpha=1$, $\beta=0$, $\phi=0$, $K_i$ very large & Attachment proportional to degree only. \\
Spatial PA & $\alpha\approx 1$, $\beta=0$, $\phi>0$, $K_i$ very large & Distance deters long links while preserving hub growth. \\
Capacity-limited PA & $\alpha\ge 0$, $\beta>0$, $\phi=0$ & Residual capacity suppresses hub formation and can truncate tails. \\
General access-based model & $\alpha\ge 0$, $\beta\ge 0$, $\phi\ge 0$, $\kappa\ge 1$ & Attachment trades off connectivity, residual capacity, and dyadic cost. \\
\bottomrule
\end{tabular}
\end{table}

\section{Simulation Design}
\label{sec:design}

All simulations use undirected simple graphs embedded in a two-dimensional unit square. Seed nodes and arriving nodes are drawn independently from the uniform distribution on $[0,1]^2$. The default seed network is a complete graph on $m_0=5$ nodes, and we require $m_0 \ge \kappa + 1$.

The revision benchmark reported here has four parts:
\begin{enumerate}[leftmargin=1.5em]
    \item four headline scenarios comparing BA, capacity-only, spatial-only, and the full generalized model;
    \item one-factor comparative statics varying $\alpha$, $\beta$, $\phi$, $\kappa$, and $K$ around a baseline;
    \item a compact $\beta$--$K$ interaction grid for heatmaps.
    \item a heterogeneous-capacity comparison that holds the mean capacity approximately fixed while varying the distribution of $K_i$ at birth.
\end{enumerate}

The headline suite uses 16 replications at $N\in\{200,1000\}$; the one-factor and $\beta$--$K$ suites use 12 replications at $N=1000$; and the heterogeneous-capacity comparison uses 16 replications at $N=1000$. The heterogeneous extension compares constant capacity $K=16$, uniform capacities $K_i\sim U[8,24]$, and lognormal capacities with mean approximately $16$. These additions do not exhaust the full design space, but they move the evidence beyond the earlier provisional benchmark. In the executed revision benchmark, no run terminated early because the feasible set became empty, but the framework records such events so that boundary conditions can be studied in future work.

\begin{table}[htbp]
\centering
\caption{Headline scenarios used throughout the paper.}
\label{tab:headline-scenarios}
\begin{tabular}{lccccc}
\toprule
Scenario & $\alpha$ & $\beta$ & $\phi$ & $\kappa$ & $K$ \\
\midrule
BA benchmark & 1.0 & 0.0 & 0.0 & 2 & very large \\
Capacity only & 1.0 & 1.0 & 0.0 & 2 & 16 \\
Spatial only & 1.0 & 0.0 & 1.0 & 2 & very large \\
General model & 1.0 & 1.0 & 1.0 & 2 & 16 \\
\bottomrule
\end{tabular}
\end{table}

For each run we record final node and edge counts, early-stop status, mean and maximum degree, degree Gini, share of nodes at capacity, connected components, size of the largest connected component, clustering, path length, diameter where feasible, assortativity, edge-length statistics, degree centralization, degree-one share, high-degree share, and tree-likeness indicators. We also save the full degree sequence of every replication.

Degree tails are assessed with more than log-log visual inspection. For each scenario we compute the empirical CCDF, choose $k_{\min}$ by Kolmogorov--Smirnov minimization for a discrete power-law tail, estimate the power-law exponent by maximum likelihood, and compare power-law, exponential, and lognormal alternatives \citep{clauset2009power}.

\section{Results}
\label{sec:results}

\subsection{Headline scenarios}

Figure~\ref{fig:headline-networks} shows one realization from each headline scenario. The BA benchmark produces the familiar hub-dominated structure. Capacity alone suppresses extreme hubs. Spatial friction alone localizes links while retaining strong hubs. The full model combines both effects: large hubs are capped and realized links remain geographically short. These layouts are illustrative; the main inferential weight comes from the comparative metrics and tail fits that follow.

\begin{figure}[htbp]
\centering
\includegraphics[width=0.92\textwidth]{figure_01_example_networks.pdf}
\caption{Final spatial network realizations for the four headline scenarios, showing how capacity and distance deterrence alter hub formation and local structure.}
\label{fig:headline-networks}
\end{figure}

The corresponding outcomes at $N=1000$ appear in Table~\ref{tab:headline-results}. Capacity is the dominant force truncating the upper tail. The mean maximum degree falls from $82.31$ in the BA benchmark to $16.00$ when $K=16$, and degree Gini falls from $0.388$ to about $0.335$. Spatial friction primarily shortens links: mean edge length falls from $0.521$ in the BA benchmark to $0.351$ in the spatial-only case. The full generalized model combines the short links of the spatial case with the capped hubs of the capacity-constrained case.

\begin{table}[htbp]
\centering
\caption{Headline outcomes at $N=1000$ (means across 16 replications).}
\label{tab:headline-results}
\begin{tabular}{lccccc}
\toprule
Scenario & Max degree & Degree Gini & Clustering & Mean edge length & Preferred tail \\
\midrule
BA benchmark & 82.31 & 0.388 & 0.028 & 0.521 & Power law \\
Capacity only & 16.00 & 0.335 & 0.007 & 0.521 & Exponential \\
Spatial only & 85.50 & 0.389 & 0.035 & 0.351 & Power law \\
General model & 16.00 & 0.336 & 0.009 & 0.349 & Exponential \\
\bottomrule
\end{tabular}
\end{table}

These results are important because the mechanisms do not merely make the same network ``look different.'' They do different structural work. Capacity suppresses hub concentration; distance localizes links and raises local closure; their combination yields a localized truncated-tail regime rather than just a more complicated BA network.

\subsection{Comparative statics}

Distance deterrence produces one of the clearest patterns in the benchmark. As $\phi$ rises from $0$ to $2$, mean edge length falls from $0.526$ to $0.157$, a reduction of roughly $70\%$, while clustering rises from $0.0075$ to $0.0721$, nearly a tenfold increase. In the explored envelope, the network remains connected throughout and no run terminates early, so stronger spatial friction bends growth locally without fragmenting the giant component.

\begin{figure}[htbp]
\centering
\includegraphics[width=0.72\textwidth]{figure_03_mean_edge_length_vs_phi.pdf}
\caption{Mean edge length as a function of the distance-deterrence parameter $\phi$, holding the other parameters at the baseline sensitivity setting.}
\label{fig:phi-edge}
\end{figure}

\begin{figure}[htbp]
\centering
\includegraphics[width=0.92\textwidth]{figure_07_connectedness_vs_phi.pdf}
\caption{Largest-component share and average path length on the largest connected component as $\phi$ varies. In the executed benchmark, the largest component always contains the full network.}
\label{fig:phi-connectedness}
\end{figure}

Multi-link entry changes the network in a different way. Moving from $\kappa=1$ to $\kappa=4$ raises average clustering from $0.00035$ to $0.0151$, raises the cyclomatic number from $6$ to $2991$, and reduces average shortest-path length from $8.19$ to $3.58$. Larger $\kappa$ therefore pushes growth away from a near-tree and toward a meshed, redundant regime.

\begin{figure}[htbp]
\centering
\includegraphics[width=0.72\textwidth]{figure_05_clustering_vs_kappa.pdf}
\caption{Average clustering coefficient as a function of the number of initial attachments $\kappa$.}
\label{fig:kappa-clustering}
\end{figure}

Capacity and saturation interact in a more subtle manner than a simple ``more $\beta$ means more saturation'' slogan would suggest. At fixed low $K$, increasing $\beta$ often \emph{reduces} the observed share of nodes exactly at capacity because the kernel diverts attachment away from nearly full nodes before they bind. For example, when $K=8$, increasing $\beta$ from $0$ to $2$ lowers the degree Gini from $0.300$ to $0.246$ and lowers the share at capacity from $0.174$ to $0.024$. Strong saturation pressure therefore changes growth by steering links, not only by enforcing hard ceilings.

\begin{figure}[htbp]
\centering
\includegraphics[width=0.72\textwidth]{figure_04_share_at_capacity_heatmap.pdf}
\caption{Share of nodes at capacity across the $\beta$--$K$ design, illustrating how stronger saturation and lower capacity change the incidence of binding constraints.}
\label{fig:capacity-saturation}
\end{figure}

\begin{figure}[htbp]
\centering
\begin{subfigure}[t]{0.48\textwidth}
    \centering
    \includegraphics[width=\textwidth]{figure_06a_max_degree_heatmap.pdf}
    \caption{Maximum degree.}
\end{subfigure}
\hfill
\begin{subfigure}[t]{0.48\textwidth}
    \centering
    \includegraphics[width=\textwidth]{figure_06b_degree_gini_heatmap.pdf}
    \caption{Degree Gini.}
\end{subfigure}
\caption{Heatmaps of maximum degree and degree inequality over the $\beta$--$K$ design. The strongest truncation occurs when low capacity and high saturation pressure operate jointly.}
\label{fig:capacity-heatmaps}
\end{figure}

The one-factor sensitivity panels reinforce a regime interpretation. Increasing $\alpha$ in the presence of a binding ceiling raises degree inequality more than it raises the maximum degree, because the ceiling already limits the largest hubs. Increasing $K$ from $8$ to effectively unbounded values restores BA-like hub growth: the mean maximum degree rises from $8.00$ to $75.00$, degree Gini rises from $0.277$ to $0.389$, and the share at capacity falls from $0.090$ to zero.

\begin{figure}[htbp]
\centering
\includegraphics[width=0.95\textwidth]{figure_08_sensitivity_panels.pdf}
\caption{One-factor sensitivity panels for $\alpha$, $\beta$, $\phi$, $\kappa$, and $K$ around the baseline specification.}
\label{fig:sensitivity-panels}
\end{figure}

\subsection{Heterogeneous capacities}

Capacity heterogeneity matters even when the mean capacity is held roughly fixed. Table~\ref{tab:heterogeneous-capacity} compares the general model under constant capacity $K=16$, uniform capacities $K_i\sim U[8,24]$, and lognormal capacities with mean approximately $16$. Relative to the constant case, heterogeneous capacities partially restore hub formation because some nodes enter with substantially more slack. The mean maximum degree rises from $16.00$ under constant capacity to $22.94$ under the uniform distribution and to $48.69$ under the lognormal distribution, while degree Gini rises from $0.337$ to $0.354$ and $0.382$, respectively. The share of nodes observed exactly at capacity falls slightly, because capacity heterogeneity diffuses binding constraints across nodes rather than forcing many nodes onto the same ceiling.

\begin{table}[htbp]
\centering
\caption{General-model comparison under homogeneous and heterogeneous capacity specifications at $N=1000$ (means across 16 replications).}
\label{tab:heterogeneous-capacity}
\begin{tabular}{lcccc}
\toprule
Capacity specification & Max degree & Degree Gini & Share at capacity & Clustering \\
\midrule
Constant $K=16$ & 16.00 & 0.337 & 0.0039 & 0.0096 \\
Uniform $U[8,24]$ & 22.94 & 0.354 & 0.0019 & 0.0121 \\
Lognormal, mean $\approx 16$ & 48.69 & 0.382 & 0.0019 & 0.0197 \\
\bottomrule
\end{tabular}
\end{table}

\subsection{Tail behavior and regime structure}

Figure~\ref{fig:ccdf} compares the empirical complementary cumulative degree distributions for the headline scenarios. The visual impression of a truncated upper tail under capacity constraints is confirmed by the tail fitting, but the BA result is more size-sensitive than in the earlier draft. The capacity-only and full generalized models prefer exponential tails at both $N=200$ and $N=1000$. The spatial-only model is mixed: exponential at $N=200$, but power-law at $N=1000$. The BA benchmark also changes with size: at $N=200$ the preferred fit is exponential, whereas at $N=1000$ it is power law with estimated tail exponent about $2.80$. The BA limit therefore reproduces the familiar heavy-tailed pattern only at the larger executed size, which is more cautious and more consistent with the finite-size concerns raised in the scale-free critique literature. In the benchmark, finite capacity remains the main force pushing the network away from BA-like heavy tails.

\begin{figure}[htbp]
\centering
\includegraphics[width=0.78\textwidth]{figure_02_degree_ccdf.pdf}
\caption{Empirical complementary cumulative degree distributions for the headline scenarios on log--log axes, using pooled degree sequences from the largest executed network size.}
\label{fig:ccdf}
\end{figure}

Taken together, the simulations identify four broad behavioral domains:
\begin{enumerate}[leftmargin=1.5em]
    \item \textbf{Hub-dominated BA-like regime}: high $K$, low $\phi$, and moderate $\kappa$ can recover familiar heavy-tailed outcomes at sufficiently large $N$.
    \item \textbf{Localized hub regime}: high $K$ with positive $\phi$ preserves hubs but shortens realized edges and raises clustering.
    \item \textbf{Saturated truncated-tail regime}: finite $K$ suppresses hub growth and degree inequality, often favoring exponential tails.
    \item \textbf{Meshed regime}: larger $\kappa$ yields more cycles, more redundancy, and shorter path lengths.
\end{enumerate}

This regime map is what the equation alone cannot deliver. The kernel tells us the direction of the forces; the simulations reveal their relative magnitude, interactions, and boundary conditions.

\section{Discussion}
\label{sec:discussion}

The integrated value of the model lies in unifying several known mechanisms around one interpretable objective. The same access-based kernel recovers BA when constraints are weak, but it also shows how capacity and space bend growth away from pure cumulative advantage. This makes the model a true generalization rather than just another attachment variant. At the same time, the present paper should be read honestly as a reduced-form access model: Equation~\ref{eq:kernel} is a separable approximation to the marginal-access ranking in Equation~\ref{eq:deltaA}, not a fully solved closed-form derivation from a complete accessibility model.

Several substantive implications follow. First, capacity matters both as a ceiling and as a steering mechanism that changes attachment probabilities before nodes saturate. Second, distance deterrence mainly localizes network growth and increases clustering in the executed benchmark, rather than fragmenting the giant component. Third, multi-link entry is crucial for moving from dendritic to meshed forms. Fourth, capacity heterogeneity partially reopens the upper tail even when average capacity is held roughly constant. These become falsifiable empirical hypotheses: stronger spatial deterrence should shorten realized links and raise clustering; stronger capacity limits should reduce degree inequality and favor exponential cutoffs; larger $\kappa$ should increase redundancy and reduce tree-likeness; and greater dispersion in node capacities should restore hub inequality relative to a homogeneous-capacity system with the same mean.

The results also clarify what the current paper does \emph{not} yet establish. We do not claim that the present benchmark proves universal fragmentation thresholds or delivers a calibrated empirical account of observed transport-network growth. Nor do we claim that the BA limit is uniformly power-law at finite size: in the executed benchmark it only clearly favors a power-law tail at $N=1000$, not at $N=200$. Instead, the paper should be read as a theory-plus-simulation contribution that sets up a broader empirical program. A natural next step is to compare observed link additions against the generalized kernel, BA, and spatial preferential attachment using both topological and spatial fit criteria.

\section{Conclusion}

This paper proposes a generalized preferential attachment model that incorporates connectivity, capacity, dyadic cost, and multi-link entry within a single access-based framework. The model nests BA as a special case, but it also produces distinct structural regimes that BA alone cannot capture. In the revision benchmark, capacity truncates the upper tail and reduces degree inequality, space shortens links and raises local clustering, multi-link entry converts tree-like growth into meshed growth, and capacity heterogeneity partially restores hub concentration. The combined model yields a localized truncated-tail regime that is especially plausible for intentionally built infrastructure systems.

The paper's contribution is therefore twofold: conceptual and computational. Conceptually, it places access rather than abstract popularity at the center of network growth. Computationally, it maps the behavioral consequences of combining capacity, cost, and multi-link formation in one model. For a paper aimed at \emph{Network Science}, that synthesis and regime mapping are the core contribution. The empirical calibration problem remains open, but the model now provides a clear platform for it.

\bibliographystyle{plainnat}
\bibliography{saturation}

\end{document}
